\begin{abstract}
\emergencystretch=2em
Affective active inference formalizes valence as inference over expected policy precision: an internal estimate of how well an agent's generative model supports action. In social interaction, however, model fitness is not only global. An agent may predict one partner reliably, remain uncertain about another, and be misled by a third. We extend this formulation by factorizing affective precision across partners in a repeated trust-game simulation. The agent maintains separate predictive models for each partner; after each interaction, partner-specific prediction error updates a local precision estimate, which modulates confidence in future policies involving that partner.

We evaluate stable interaction, partner choice, and abrupt betrayal regimes. Partner-local precision changes how social beliefs are acted upon, shaping which partners the agent returns to, how decisively it acts, and how it reallocates after surprising outcomes. These effects are strongest in decision-ambiguous regimes, where multiple actions remain viable. They can also emerge even when agents with and without precision modulation maintain similar beliefs about their partners, indicating a deployment rather than knowledge difference. Because this affective signal tracks predictive reliability more than partner reward, its behavioral influence supports a model-fitness interpretation. Under abrupt betrayal, however, surprise-driven precision updates can over-sharpen decisions, locking the agent into outdated assumptions and reducing payoffs after the behavioral switch. These results suggest that extending affective inference to multi-agent systems requires localizing precision over social models, providing a computational account of how implicit metacognitive estimates of model fitness can both support and misdirect social action.
\keywords{Active inference \and Affect \and Trust games}
\end{abstract}
